\section{Discussion}
\label{sec:discussion}

Three patterns recur across studies that were designed independently and
address unrelated questions. Their recurrence is itself informative about how
this class of model should be studied going forward.

\paragraph{A reasoned hypothesis still requires a control arm.} Three separate
studies advanced a specific, mechanistically plausible hypothesis and then
tested it directly: gradient-norm roughness as the cause of Euler's
convergence floor (Section~\ref{sec:track-b}), SiLU's representational limit
as the cause of the pendulum's early failure (Section~\ref{sec:phase2}), and
the predicted inapplicability of the vanishing-field gate to non-mechanistic
data (Section~\ref{sec:track-e}). Each hypothesis had a clear argument behind it
before any run was launched, and each was overturned by the run itself. The
result reported in every one of these three cases is the outcome of the
control arm, and the original argument survives only as the motivation for
running it.

\paragraph{A single seed and a single train/test split account for results
that first looked like properties of the solver or the method.} The
stiffness-map failures at $\mu\in\{0.5,1.0,2.0\}$ (Section~\ref{sec:track-d})
and the epidemic-fit extrapolation runaway (Section~\ref{sec:track-e}) both
appeared, on first inspection, to indict a solver or a modeling approach.
Retraining with additional seeds, in the stiffness-map case, and shifting the
train/test split by fifteen days, in the epidemic case, resolved both into an
artifact of one specific run's random initialization or one specific window
boundary. Neither observation survives as a claim about solver order or about
KAN-ODE's capacity once the seed or the split is varied.

\paragraph{Measuring a quantity before training is more informative than
tuning after a failed run.} The SIR divergence (Section~\ref{sec:phase2}), the
real-epidemic-data fit (Section~\ref{sec:track-e}), and the stiffness map's
spurious equilibrium (Section~\ref{sec:track-d}) were each root-caused by
evaluating the model's own initial field, or its trained field at a candidate
equilibrium, and comparing that number against the system's true scale. In
every case this produced a specific, falsifiable prediction -- a rescaling
factor, a fix that should or should not transfer, an equilibrium location --
that a subsequent short run confirmed or refuted directly, in contrast to the
open-ended search a hyperparameter sweep would have required.

\FloatBarrier
